% Options for packages loaded elsewhere
\PassOptionsToPackage{unicode}{hyperref}
\PassOptionsToPackage{hyphens}{url}
\documentclass[
  a4paper,
]{article}
\usepackage{xcolor}
\usepackage[margin=25mm]{geometry}
\usepackage{amsmath,amssymb}
\setcounter{secnumdepth}{-\maxdimen} % remove section numbering
\usepackage{iftex}
\ifPDFTeX
  \usepackage[T1]{fontenc}
  \usepackage[utf8]{inputenc}
  \usepackage{textcomp} % provide euro and other symbols
\else % if luatex or xetex
  \usepackage{unicode-math} % this also loads fontspec
  \defaultfontfeatures{Scale=MatchLowercase}
  \defaultfontfeatures[\rmfamily]{Ligatures=TeX,Scale=1}
\fi
\usepackage[]{lmodern}
\ifPDFTeX\else
  % xetex/luatex font selection
\fi
% Use upquote if available, for straight quotes in verbatim environments
\IfFileExists{upquote.sty}{\usepackage{upquote}}{}
\IfFileExists{microtype.sty}{% use microtype if available
  \usepackage[]{microtype}
  \UseMicrotypeSet[protrusion]{basicmath} % disable protrusion for tt fonts
}{}
\makeatletter
\@ifundefined{KOMAClassName}{% if non-KOMA class
  \IfFileExists{parskip.sty}{%
    \usepackage{parskip}
  }{% else
    \setlength{\parindent}{0pt}
    \setlength{\parskip}{6pt plus 2pt minus 1pt}}
}{% if KOMA class
  \KOMAoptions{parskip=half}}
\makeatother
\usepackage{longtable,booktabs,array}
\newcounter{none} % for unnumbered tables
\usepackage{calc} % for calculating minipage widths
% Correct order of tables after \paragraph or \subparagraph
\usepackage{etoolbox}
\makeatletter
\patchcmd\longtable{\par}{\if@noskipsec\mbox{}\fi\par}{}{}
\makeatother
% Allow footnotes in longtable head/foot
\IfFileExists{footnotehyper.sty}{\usepackage{footnotehyper}}{\usepackage{footnote}}
\makesavenoteenv{longtable}
\usepackage{graphicx}
\makeatletter
\newsavebox\pandoc@box
\newcommand*\pandocbounded[1]{% scales image to fit in text height/width
  \sbox\pandoc@box{#1}%
  \Gscale@div\@tempa{\textheight}{\dimexpr\ht\pandoc@box+\dp\pandoc@box\relax}%
  \Gscale@div\@tempb{\linewidth}{\wd\pandoc@box}%
  \ifdim\@tempb\p@<\@tempa\p@\let\@tempa\@tempb\fi% select the smaller of both
  \ifdim\@tempa\p@<\p@\scalebox{\@tempa}{\usebox\pandoc@box}%
  \else\usebox{\pandoc@box}%
  \fi%
}
% Set default figure placement to htbp
\def\fps@figure{htbp}
\makeatother
\setlength{\emergencystretch}{3em} % prevent overfull lines
\providecommand{\tightlist}{%
  \setlength{\itemsep}{0pt}\setlength{\parskip}{0pt}}
\usepackage{bookmark}
\IfFileExists{xurl.sty}{\usepackage{xurl}}{} % add URL line breaks if available
\urlstyle{same}
\hypersetup{
  hidelinks,
  pdfcreator={LaTeX via pandoc}}

\author{}
\date{}

\begin{document}

\section{Linear Upper Bound on Generator Rank and Three-Direction
Saturation of Spatial-Direction Readout in N-Body Complete Pairwise
Relational
Waves}\label{linear-upper-bound-on-generator-rank-and-three-direction-saturation-of-spatial-direction-readout-in-n-body-complete-pairwise-relational-waves}

\subsection{Relational waves grow quadratically, rotation modes grow
linearly, and uniquely readable spatial directions stop at
three}\label{relational-waves-grow-quadratically-rotation-modes-grow-linearly-and-uniquely-readable-spatial-directions-stop-at-three}

\textbf{Author:} Noriaki Kihara\\
\textbf{ORCID:} 0009-0004-6753-4020\\
\textbf{Date:} July 21, 2026\\
\textbf{Version DOI:} 10.5281/zenodo.21465899\\
\textbf{Concept DOI:} 10.5281/zenodo.21465898\\
\textbf{Position:} Dimensional Generation Structure series, Paper 3, v1

\begin{center}\rule{0.5\linewidth}{0.5pt}\end{center}

\subsection{Abstract}\label{abstract}

Previous work extended the one-dimensional phase-relation readout of the
closed AB two-body phase system {[}2,3{]} to complete pairwise
relational waves of ABC three-body and ABCD four-body systems, showing
that ABC yields one rotation plane with an invariant Z normal and that
ABCD yields three rotation planes read as XYZ {[}4{]}. This paper
extends the construction to general \(N\) bodies and shows, by theory
and numerical experiment, that three quantities grow at different
orders: the number of internal relational waves, the number of rotation
modes, and the number of uniquely readable spatial directions.

The number of complete pairwise relational waves grows as
\(M_N=\binom{N}{2}=O(N^2)\). The phase-difference sine antisymmetric
generator \(\widetilde K\), however, decomposes exactly into a sum of
\(N\) antisymmetric matrices of rank at most 2 built from the rows of
the incidence matrix,

\[
\widetilde K=\sum_{k=1}^{N}\left(c_k s_k^{\mathsf T}-s_k c_k^{\mathsf T}\right),
\]

so that

\[
\operatorname{rank}K\le 2\min\!\left(N,\left\lfloor\frac{M_N}{2}\right\rfloor\right)=O(N)
\]

holds for all \(N\). For \(3\le N\le12\) we further prove that equality
holds for almost all initial phases, combining exact rational witness
configurations (tan half-angle parametrization, fraction arithmetic)
with the measure-zero property of zero sets of real-analytic functions
{[}9{]}. Finally, for the direction readout that reads a normal from the
plane spanned by two independent relational directions, the
normal-candidate space has dimension \(d-2\) for a display dimension
\(d\), so the normal is unique up to sign only for \(d=3\). Under the
present two-body phase-relation readout rule, therefore, even though
internal relational waves and internal rotation modes increase with
\(N\), the uniquely readable spatial directions saturate at three.

In numerical experiments with 32 trials of 720 steps each for \(N=3\) to
\(9\), the maximum squared-closure error was \(1.82\times10^{-13}\), the
maximum absolute-square-sum drift \(2.25\times10^{-13}\), and the
maximum relabeling-covariance error \(1.12\times10^{-13}\); the residual
of the vertex decomposition identity was at most \(5.00\times10^{-16}\).
The generic rank law held without exception in 2560 trials for \(N=3\)
to \(12\). The multi-dimensional null space reappearing for \(N\ge6\)
was confirmed directly to be a residual subspace whose projector and
squared quantities are unique while no unique internal basis exists. No
sequential state renormalization, no observation damping, and no
absolute background axes were used.

\begin{center}\rule{0.5\linewidth}{0.5pt}\end{center}

\subsection{0. Conclusions}\label{conclusions}

In the \(N\)-body complete pairwise relational-wave system, the
following three layers grow at different orders.

\[
\boxed{
M_N=\binom{N}{2}=O(N^2)
\quad\longrightarrow\quad
r\le N=O(N)
\quad\longrightarrow\quad
D_{\mathrm{unique}}=3=O(1)
}
\]

The first layer is the number of internal relational waves, which grows
quadratically.

The second layer is the number of internal rotation modes
\(r=\frac12\operatorname{rank}K\). By the vertex decomposition theorem,
\(\operatorname{rank}K\le2\min(N,\lfloor M_N/2\rfloor)\) holds for all
\(N\), so this layer grows at most linearly. For \(3\le N\le12\) we
prove rigorously that this bound is attained with equality in general
position.

The third layer is the number of uniquely readable spatial directions.
The normal-candidate space of a two-relation plane has dimension \(d-2\)
for display dimension \(d\), and sign-free uniqueness holds only for
\(d=3\). In general \(D_{\mathrm{unique}}\le3\); in general position,
where two independent and distinguishable relational directions exist,
\(D_{\mathrm{unique}}=3\). Equality fails for \(N=2\) and for degenerate
phases at which the generator vanishes. Therefore,

\[
\boxed{
\text{under the present two-body phase-relation readout rule,}
\quad
\text{the uniquely readable spatial directions saturate at three.}
}
\]

For \(N\ge6\) the null space reappears; in general position its
dimension grows as

\[
\dim\ker K=M_N-2N=\frac{N(N-5)}{2}=O(N^2).
\]

This null space is a residual subspace whose projector and squared
quantities are uniquely readable while its individual internal
directions are not. Added relational waves do not disappear: they are
retained in this basis-free internal subspace and as internal phase
modes that are not adopted as spatial directions.

These implement the adoption conditions of Axiom 16 of Basic Axiom
System v5, namely distinguishability by rotation frequencies and
uniqueness of the plane normal {[}1{]}.

\begin{center}\rule{0.5\linewidth}{0.5pt}\end{center}

\subsection{1. Research Question}\label{research-question}

\subsubsection{1.1 The question left open by the preceding
paper}\label{the-question-left-open-by-the-preceding-paper}

The preceding ABCD experiment {[}4{]} showed that six complete pairwise
relational waves decompose into three rotation planes, giving a
structure readable as three XYZ directions. It left open the following
question.

\begin{quote}
When relational waves are increased to five or more bodies, what is the
quantity that increases internally, and where does the quantity that can
be uniquely read as spatial directions stop?
\end{quote}

Reference {[}4{]} offered the interpretation that what increases beyond
four bodies are internal relations without uniquely assignable
directions, but this was outside its numerical scope and was not
established as a theorem. The present paper promotes that interpretation
to three theorems: the rank upper bound via vertex decomposition, the
generic-position equality, and normal uniqueness.

\subsubsection{1.2 Rules on the direction of
causation}\label{rules-on-the-direction-of-causation}

The causal direction of this paper is

\[
\boxed{
\text{relational waves and squared closure}
\longrightarrow
\text{readout maps}
\longrightarrow
\text{uniquely reconstructible directions}
\longrightarrow
\text{spatial display}.
}
\]

Space does not exist at the outset. Accordingly, this paper does not do
any of the following.

\begin{enumerate}
\def\labelenumi{\arabic{enumi}.}
\tightlist
\item
  It does not reinterpret a relational wave \(X_{ij}\) as a distance
  \(d_{ij}\) in a background space.
\item
  It does not use the embedding rank into a pre-existing Euclidean space
  as a criterion for the generation of space.
\item
  It does not identify the number of rotation planes, the generator
  rank, or the null-space dimension with the number of spatial axes.
\item
  It does not add constraints of external geometric origin at five or
  more bodies. All relational waves participate in the same single
  squared closure.
\end{enumerate}

What this paper judges is not the size of the state space but the
uniqueness of the readout map.

\subsubsection{1.3 What this paper
discriminates}\label{what-this-paper-discriminates}

\begin{enumerate}
\def\labelenumi{\arabic{enumi}.}
\tightlist
\item
  Can the all-\(N\) upper bound on generator rank be derived from inside
  the model?
\item
  Does the bound become an equality in general position, and can this be
  proven rigorously?
\item
  In which display dimensions is the normal readout unique?
\item
  As what kind of quantity can the null space reappearing at \(N\ge6\)
  be read?
\item
  Are the conserved quantities and relabeling covariance maintained
  after the \(N\) extension?
\end{enumerate}

\begin{center}\rule{0.5\linewidth}{0.5pt}\end{center}

\subsection{2. Classification of Claims}\label{classification-of-claims}

This paper distinguishes definitions, mathematical consequences,
numerical facts, spatial-direction interpretations, and items left
unfixed. The axiomatic basis of anonymity, squared closure, and complete
pairwise relational waves is Basic Axiom System v5 {[}1{]}.

{\def\LTcaptype{none} % do not increment counter
\begin{longtable}[]{@{}
  >{\raggedright\arraybackslash}p{(\linewidth - 4\tabcolsep) * \real{0.3333}}
  >{\raggedright\arraybackslash}p{(\linewidth - 4\tabcolsep) * \real{0.3333}}
  >{\raggedright\arraybackslash}p{(\linewidth - 4\tabcolsep) * \real{0.3333}}@{}}
\toprule\noalign{}
\begin{minipage}[b]{\linewidth}\raggedright
Item
\end{minipage} & \begin{minipage}[b]{\linewidth}\raggedright
Classification
\end{minipage} & \begin{minipage}[b]{\linewidth}\raggedright
Position in this paper
\end{minipage} \\
\midrule\noalign{}
\endhead
\bottomrule\noalign{}
\endlastfoot
Complete pairwise relational waves as independent state waves & Model
definition & Implementation of Axiom 16 of {[}1{]}; inherited from
{[}4{]} \\
\(\sum_e X_e^2+(iR)^2=0\) & Closure condition & Computed directly at
every step \\
Phase-difference sine generator coupling only endpoint-sharing relations
& Working hypothesis & Same as {[}4{]}; not derived from the axioms \\
Vertex decomposition
\(\widetilde K=\sum_k(c_ks_k^{\mathsf T}-s_kc_k^{\mathsf T})\) &
Mathematical consequence & Theorem 1; numerical residual
\(5.00\times10^{-16}\) \\
\(\operatorname{rank}K\le2\min(N,\lfloor M/2\rfloor)\) (all \(N\)) &
Mathematical consequence & Corollary of Theorem 1 \\
Generic equality (\(3\le N\le12\)) & Mathematical consequence
(computer-assisted proof) & Theorem 2; exact rational witnesses +
measure-zero argument \\
Generic equality (all \(N\)) & New hypothesis & General witness
construction unproven \\
Normal uniqueness (only \(d=3\)) & Mathematical consequence & Theorem
3 \\
Three-direction saturation (\(D_{\mathrm{unique}}\le3\), \(=3\) in
general position) & Consequence of Axiom 16 and Theorem 3 & Section 7;
does not follow from generator rank alone \\
Results of Experiments A--E & Numerical facts & Section 8 \\
Reading three distinguishable rotation planes as XYZ & Spatial-direction
interpretation & Inherited from {[}4{]} \\
Reading residues as internal phase modes and residual subspaces &
Spatial-direction interpretation & Consistent with the residue clause of
Axiom 16 {[}1{]} \\
Physical axes corresponding to residual directions & Out of scope & Not
identified in this paper \\
\end{longtable}
}

\begin{center}\rule{0.5\linewidth}{0.5pt}\end{center}

\subsection{3. The N-Body Complete Pairwise Relational-Wave
Model}\label{the-n-body-complete-pairwise-relational-wave-model}

\subsubsection{3.1 Relation set and relational
waves}\label{relation-set-and-relational-waves}

For \(N\) anonymous individuals, the set of all pairwise relations is

\[
\mathcal E_N
=
\bigl\{\{i,j\}\mid 1\le i<j\le N\bigr\},
\qquad
M=\lvert\mathcal E_N\rvert=\binom{N}{2}.
\]

To each relation \(e\) we assign a complex relational wave
\(X_e=a_e+ib_e\) and write the full state as \(X\in\mathbb C^M\).
Relational waves are not observation channels; they are independent
physical state components not derived from component waves {[}1{]}.

The closure condition is

\[
\sum_{e\in\mathcal E_N}X_e^2+(iR)^2=0,
\qquad
\text{transposed form}\quad
q(X)=X^{\mathsf T}X=R^2.
\]

As an independent comparison quantity we also record the absolute-square
sum \(H(X)=X^\dagger X\).

\subsubsection{3.2 Adjacency and
generator}\label{adjacency-and-generator}

Two relations \(e,f\) are adjacent exactly when they share one endpoint:

\[
A_{ef}
=
\begin{cases}
1,& e\ne f\ \text{and}\ e\cap f\ne\varnothing,\\
0,& \text{otherwise}.
\end{cases}
\]

This is the adjacency matrix of the line graph of the complete graph
\(K_N\) {[}5{]}.

From the initial phases \(\theta_e=\arg X_e(0)\), the unnormalized
generator is

\[
\widetilde K_{ef}
=
A_{ef}\sin(\theta_f-\theta_e).
\]

It satisfies \(\widetilde K^{\mathsf T}=-\widetilde K\); when nonzero it
is divided by its spectral norm to give \(K\). This phase-difference
sine coupling rule is the same conservative working hypothesis as in
{[}4{]}. The dynamics of this paper uses a fixed generator: the
generator is constructed once from the initial phases and kept fixed
thereafter.

The one-step update is the real orthogonal matrix given by the Cayley
transform

\[
U=(I-\gamma K)^{-1}(I+\gamma K),
\qquad
U^{\mathsf T}U=I
\]

{[}7{]}. Orthogonality preserves \(q(X)\) and \(H(X)\) without
sequential renormalization.

\subsubsection{3.3 Incidence matrix}\label{incidence-matrix}

Let \(B\in\mathbb R^{N\times M}\) be the unsigned incidence matrix with
individuals as rows and relations as columns; the column of relation
\(e=\{i,j\}\) has ones exactly in rows \(i,j\). It relates to the
line-graph adjacency by

\[
A=B^{\mathsf T}B-2I_M.
\]

\begin{center}\rule{0.5\linewidth}{0.5pt}\end{center}

\subsection{4. Vertex Decomposition
Theorem}\label{vertex-decomposition-theorem}

\subsubsection{4.1 Theorem}\label{theorem}

\begin{quote}
\textbf{Theorem 1 (vertex decomposition)}

Let \(b_k\) be the \(k\)-th row of \(B\) and set

\[
c_k=(B_{ke}\cos\theta_e)_{e\in\mathcal E_N},
\qquad
s_k=(B_{ke}\sin\theta_e)_{e\in\mathcal E_N}.
\]

Then

\[
\widetilde K=\sum_{k=1}^{N}\left(c_k s_k^{\mathsf T}-s_k c_k^{\mathsf T}\right)
\]

holds exactly.
\end{quote}

\textbf{Proof.} By entrywise computation,

\[
\left(c_k s_k^{\mathsf T}-s_k c_k^{\mathsf T}\right)_{ef}
=
B_{ke}B_{kf}\left(\cos\theta_e\sin\theta_f-\sin\theta_e\cos\theta_f\right)
=
B_{ke}B_{kf}\sin(\theta_f-\theta_e).
\]

Summing over \(k\), the coefficient is
\(\sum_kB_{ke}B_{kf}=(B^{\mathsf T}B)_{ef}\). For \(e=f\), \(\sin0=0\)
kills the diagonal. For \(e\ne f\), two distinct edges share at most one
endpoint, so \((B^{\mathsf T}B)_{ef}=A_{ef}\in\{0,1\}\). Every entry
therefore equals \(A_{ef}\sin(\theta_f-\theta_e)=\widetilde K_{ef}\).
\(\blacksquare\)

In the numerical experiments the residual of this identity was at most
\(5.00\times10^{-16}\) (machine precision) over all configurations and
trials.

\subsubsection{4.2 Linear upper bound on the
rank}\label{linear-upper-bound-on-the-rank}

\begin{quote}
\textbf{Corollary 1 (rank bound)}

For all \(N\) and all initial phases,

\[
\operatorname{rank}K
\le
2\min\!\left(N,\left\lfloor\frac{M}{2}\right\rfloor\right).
\]
\end{quote}

\textbf{Proof.} Each term \(c_ks_k^{\mathsf T}-s_kc_k^{\mathsf T}\) in
Theorem 1 has rank at most 2, so the sum has rank at most \(2N\). The
bound \(\operatorname{rank}K\le M\) is trivial. A real antisymmetric
matrix decomposes under an orthogonal transformation into
two-dimensional rotation blocks and a zero block, so its rank is even
{[}6{]}. Combining the three gives the bound. \(\blacksquare\)

Corollary 1 shows, from inside the model and without introducing a
background space, that while the number of relational waves grows as
\(O(N^2)\), the dynamical rotational degrees of freedom are at most
\(O(N)\).

\textbf{Remark (why the bound is saturable).} \(K_N\) is connected and
non-bipartite for \(N\ge3\), and the real rank of the unsigned incidence
matrix is then \(N\) {[}8{]}. The \(2N\) vectors \(\{c_k,s_k\}\)
spanning the column space of \(\widetilde K\) are thus built from \(N\)
independent incidence rows, so there is no degeneration on the
coefficient side. As shown in the next section, the bound is actually
attained in general position.

\begin{center}\rule{0.5\linewidth}{0.5pt}\end{center}

\subsection{5. Generic-Position Rank
Theorem}\label{generic-position-rank-theorem}

\subsubsection{5.1 Theorem}\label{theorem-1}

\begin{quote}
\textbf{Theorem 2 (generic equality, \(3\le N\le12\))}

For each \(N\) with \(3\le N\le12\), setting
\(m=\min(N,\lfloor M/2\rfloor)\), for almost all initial phases
\(\theta\in[0,2\pi)^M\) (all except a set of Lebesgue measure zero),

\[
\operatorname{rank}K=2m.
\]
\end{quote}

\textbf{Proof.} The upper bound is Corollary 1. For the lower bound, the
entries of \(\widetilde K\) are

\[
\sin(\theta_f-\theta_e)
=
\sin\theta_f\cos\theta_e-\cos\theta_f\sin\theta_e,
\]

polynomials in the unit-circle points \((\cos\theta_e,\sin\theta_e)\).
Using the tan half-angle parameter \(t_e\),

\[
(\cos\theta_e,\sin\theta_e)
=
\left(\frac{1-t_e^2}{1+t_e^2},\ \frac{2t_e}{1+t_e^2}\right),
\]

the witness configuration \(t_e=e\) (\(e=1,\dots,M\)) makes
\(\widetilde K\) an exactly rational matrix. Computing its rank exactly
by Gaussian elimination in fraction arithmetic yields

\[
\operatorname{rank}\widetilde K=2m
\]

for all \(3\le N\le12\) (verification program and output in Section 10).
Since the computation contains no floating point, this is an existence
proof of a witness for each \(N\).

Because \(\operatorname{rank}\widetilde K=2m\) at the witness, some
\(2m\times2m\) minor \(D(\theta)\) is nonzero there. \(D\) is a
real-analytic function of \(\theta\) and is not identically zero. The
zero set of a nontrivial real-analytic function has Lebesgue measure
zero {[}9{]}. The set of phases with rank below \(2m\) is contained in
\(\{\theta\mid D(\theta)=0\}\) and hence has measure zero.
\(\blacksquare\)

\subsubsection{5.2 Generalization to all
N}\label{generalization-to-all-n}

\begin{quote}
\textbf{New hypothesis (generic equality, all \(N\))}

For all \(N\ge3\), \(\operatorname{rank}K=2\min(N,\lfloor M/2\rfloor)\)
in general position.
\end{quote}

This is proven for \(N\le12\) by Theorem 2. A general witness
construction for arbitrary \(N\) is unproven, and the claim is stated
explicitly as a new hypothesis. The same exact rational verification
extends mechanically to any individual \(N\).

\subsubsection{5.3 Degenerate
configurations}\label{degenerate-configurations}

Away from general position the rank can drop. Experiment B confirmed the
following.

\begin{enumerate}
\def\labelenumi{\arabic{enumi}.}
\tightlist
\item
  When all phases are equal, or restricted to \(\{0,\pi\}\), all
  \(\sin(\theta_f-\theta_e)=0\) and the generator vanishes (rank 0).
\item
  With two-valued phases \(\{0,0.5\}\), the rank dropped from \(10\) to
  \(8\) only at \(N=5\). \(N=5\) is the critical configuration \(M=2N\);
  this observation is a concrete example showing that the exceptional
  set of Theorem 2 is not empty.
\end{enumerate}

\begin{center}\rule{0.5\linewidth}{0.5pt}\end{center}

\subsection{6. Normal-Uniqueness
Theorem}\label{normal-uniqueness-theorem}

\subsubsection{6.1 Definition of a unique
direction}\label{definition-of-a-unique-direction}

A direction is uniquely determined when the following hold.

\begin{enumerate}
\def\labelenumi{\arabic{enumi}.}
\tightlist
\item
  The direction can be reconstructed from relational data alone.
\item
  No absolute background axis is referenced.
\item
  It transforms covariantly under permutations of individual names.
\item
  No additional rule selecting one candidate among many is required.
\end{enumerate}

Uniqueness here is first the uniqueness of a one-dimensional subspace up
to sign.

\subsubsection{6.2 Two-relation plane and normal-candidate
space}\label{two-relation-plane-and-normal-candidate-space}

Let \(u,v\in V_d\) be two linearly independent relational directions in
a real \(d\)-dimensional display space, spanning the plane
\(P=\operatorname{span}\{u,v\}\). With
\(A_P=\begin{pmatrix}u^{\mathsf T}\\ v^{\mathsf T}\end{pmatrix}\), the
normal-candidate space is

\[
\mathcal N(P)=\ker A_P=P^\perp,
\]

and its orthogonal projector

\[
P_\perp=I_d-A_P^{\mathsf T}\left(A_PA_P^{\mathsf T}\right)^{-1}A_P
\]

is determined by \(u,v\) alone, without external axis names.

\subsubsection{6.3 Theorem}\label{theorem-2}

\begin{quote}
\textbf{Theorem 3 (normal uniqueness)}

When a normal direction is reconstructed from the plane spanned by two
linearly independent relational directions, without additional
background axes or selection rules, the normal is unique up to sign only
when the display dimension is

\[
d=3.
\]
\end{quote}

\textbf{Proof.} Since \(\operatorname{rank}A_P=2\), the rank--nullity
theorem gives \(\dim\mathcal N(P)=d-2\).

For \(d=3\), \(\dim\mathcal N(P)=1\), and the projector
\(P_\perp=nn^{\mathsf T}\) built from the normalized normal \(n\) is
invariant under sign reversal, \((-n)(-n)^{\mathsf T}=nn^{\mathsf T}\);
the direction line is unique.

For \(d\ge4\), \(\dim\mathcal N(P)=d-2\ge2\). The group of orthogonal
transformations fixing \(u,v\) acts on the normal-candidate space as
\(O(d-2)\), mixing normal candidates while leaving \(P_\perp\)
unchanged. The set of unit normal candidates is \(S^{d-3}\), and after
sign identification \(\mathbb{RP}^{d-3}\) --- a continuum, in contrast
to the single point \(\mathbb{RP}^0\) at \(d=3\). Relational data do not
distinguish elements of this family, so selecting one is incompatible
with anonymity. \(\blacksquare\)

Even for \(d\ge4\), the projector \(P_\perp\) and the squared sum of
candidate components remain uniquely readable. What cannot be read
uniquely are the individual linear directions inside.

\begin{center}\rule{0.5\linewidth}{0.5pt}\end{center}

\subsection{7. Three-Layer Hierarchy and Residual
Subspace}\label{three-layer-hierarchy-and-residual-subspace}

\subsubsection{7.1 Composition of the three
layers}\label{composition-of-the-three-layers}

Applying Theorem 3 to the internal mode structure given by Theorems 1
and 2 and to the observation selection based on Axiom 16 {[}1{]}, the
\(N\)-body complete pairwise relational-wave system separates into three
layers.

{\def\LTcaptype{none} % do not increment counter
\begin{longtable}[]{@{}
  >{\raggedright\arraybackslash}p{(\linewidth - 6\tabcolsep) * \real{0.2500}}
  >{\raggedright\arraybackslash}p{(\linewidth - 6\tabcolsep) * \real{0.2500}}
  >{\raggedright\arraybackslash}p{(\linewidth - 6\tabcolsep) * \real{0.2500}}
  >{\raggedright\arraybackslash}p{(\linewidth - 6\tabcolsep) * \real{0.2500}}@{}}
\toprule\noalign{}
\begin{minipage}[b]{\linewidth}\raggedright
Layer
\end{minipage} & \begin{minipage}[b]{\linewidth}\raggedright
Quantity
\end{minipage} & \begin{minipage}[b]{\linewidth}\raggedright
Order
\end{minipage} & \begin{minipage}[b]{\linewidth}\raggedright
Basis
\end{minipage} \\
\midrule\noalign{}
\endhead
\bottomrule\noalign{}
\endlastfoot
Internal relational waves & \(M_N=\binom{N}{2}\) & \(O(N^2)\) & Model
definition \\
Internal rotation modes & \(r=\frac12\operatorname{rank}K\le N\) &
\(O(N)\) & Theorems 1, 2 \\
Uniquely readable spatial directions & \(D_{\mathrm{unique}}\le3\)
(\(=3\) in general position) & \(O(1)\) & Axiom 16 + Theorem 3 \\
\end{longtable}
}

Three-direction saturation is not a consequence of the generator rank
alone. It holds only when the observation selection of Axiom 16
(distinguishability plus normal uniqueness) is connected with Theorem 3.
Nor does \(D_{\mathrm{unique}}=3\) hold always: it holds in general
position where two independent and distinguishable relational directions
exist, and equality fails for \(N=2\) or for degenerate phases at which
the generator vanishes.

The number of rotation planes is not the number of spatial directions. A
rotation plane becomes a candidate direction when distinguishable by its
independent rotation frequency, but adoption as spatial directions is
limited to the three directions satisfying normal uniqueness. When
distinguishable rotation planes exceed that range, the excess planes are
retained as internal phase modes. This implements the adoption
conditions of Axiom 16 of Basic Axiom System v5 {[}1{]}.

\subsubsection{7.2 Growth of the residual
subspace}\label{growth-of-the-residual-subspace}

In general position, for \(N\ge5\),

\[
\dim\ker K
=
M_N-2N
=
\frac{N(N-5)}{2}
\]

(1 for \(N=3\); 0 for \(N=4,5\)). The null space reappears at \(N\ge6\)
and grows as \(O(N^2)\). For large \(N\), most of the relational-wave
space becomes a residual subspace unmoved by the generator.

The one-dimensional kernel of ABC could be read uniquely as the Z normal
{[}4{]}. The kernel for \(N\ge6\) is multi-dimensional: its orthogonal
projector and projection norm are unique, but its internal directions
carry an \(O(\dim\ker K)\) gauge freedom and cannot be selected
uniquely. The residue is therefore read as

\[
\boxed{
\text{not unique linear directions, but an internal subspace without a unique basis,}
}
\]

consistent with the residue clause of Axiom 16 {[}1{]}.

\subsubsection{7.3 What is preserved}\label{what-is-preserved}

What the \(N\) extension preserves is not an X, Y, Z with particular
names. It is the three-direction structure consisting of two independent
relational directions selected from distinguishable rotation planes, and
one normal direction uniquely determined by their plane. Under name
permutations the relational components composing the three directions
move covariantly.

\begin{center}\rule{0.5\linewidth}{0.5pt}\end{center}

\subsection{8. Numerical Experiments}\label{numerical-experiments}

\subsubsection{8.1 Common conditions}\label{common-conditions}

{\def\LTcaptype{none} % do not increment counter
\begin{longtable}[]{@{}lr@{}}
\toprule\noalign{}
Item & Value \\
\midrule\noalign{}
\endhead
\bottomrule\noalign{}
\endlastfoot
Dynamics & Fixed generator (constructed once from initial phases) \\
\(R^2\) & 1.0 \\
Imaginary seed amplitude \(s\) & 0.35 \\
Cayley coefficient \(\gamma\) & \(\tan(\pi/144)\) \\
Random seed & 20260721 (Experiments B--E: +1 to +4) \\
Tolerance & \(10^{-10}\) \\
Sequential state renormalization & None \\
Observation damping & None \\
Absolute background axes & None \\
\end{longtable}
}

The initial state is \(X(0)=\sqrt{R^2+s^2}\,u+isv\) with real unit
vectors \(u\perp v\), for which \(q(X(0))=R^2\) holds exactly {[}4{]}.

\subsubsection{8.2 Experiment A: conservation, covariance, spectral
structure for N =
3--9}\label{experiment-a-conservation-covariance-spectral-structure-for-n-39}

32 trials of 720 steps each per \(N\).

{\def\LTcaptype{none} % do not increment counter
\begin{longtable}[]{@{}rrrrrl@{}}
\toprule\noalign{}
\(N\) & \(M\) & rank \(K\) & \(\dim\ker K\) & Rotation planes & Rank
law \\
\midrule\noalign{}
\endhead
\bottomrule\noalign{}
\endlastfoot
3 & 3 & 2 & 1 & 1 & holds \\
4 & 6 & 6 & 0 & 3 & holds \\
5 & 10 & 10 & 0 & 5 & holds \\
6 & 15 & 12 & 3 & 6 & holds \\
7 & 21 & 14 & 7 & 7 & holds \\
8 & 28 & 16 & 12 & 8 & holds \\
9 & 36 & 18 & 18 & 9 & holds \\
\end{longtable}
}

Ranks and null-space dimensions agreed across all 32 trials for each
\(N\), all equal to \(2\min(N,\lfloor M/2\rfloor)\) and \(M-2N\)
(\(N\ge5\)).

{\def\LTcaptype{none} % do not increment counter
\begin{longtable}[]{@{}
  >{\raggedleft\arraybackslash}p{(\linewidth - 12\tabcolsep) * \real{0.1429}}
  >{\raggedleft\arraybackslash}p{(\linewidth - 12\tabcolsep) * \real{0.1429}}
  >{\raggedleft\arraybackslash}p{(\linewidth - 12\tabcolsep) * \real{0.1429}}
  >{\raggedleft\arraybackslash}p{(\linewidth - 12\tabcolsep) * \real{0.1429}}
  >{\raggedleft\arraybackslash}p{(\linewidth - 12\tabcolsep) * \real{0.1429}}
  >{\raggedleft\arraybackslash}p{(\linewidth - 12\tabcolsep) * \real{0.1429}}
  >{\raggedleft\arraybackslash}p{(\linewidth - 12\tabcolsep) * \real{0.1429}}@{}}
\toprule\noalign{}
\begin{minipage}[b]{\linewidth}\raggedleft
\(N\)
\end{minipage} & \begin{minipage}[b]{\linewidth}\raggedleft
Max closure error
\end{minipage} & \begin{minipage}[b]{\linewidth}\raggedleft
Max \(H\) drift
\end{minipage} & \begin{minipage}[b]{\linewidth}\raggedleft
Max orthogonality error
\end{minipage} & \begin{minipage}[b]{\linewidth}\raggedleft
Max trajectory covariance error
\end{minipage} & \begin{minipage}[b]{\linewidth}\raggedleft
Vertex decomposition residual
\end{minipage} & \begin{minipage}[b]{\linewidth}\raggedleft
Max kernel projection drift
\end{minipage} \\
\midrule\noalign{}
\endhead
\bottomrule\noalign{}
\endlastfoot
3 & \(1.44\times10^{-13}\) & \(1.78\times10^{-13}\) &
\(4.44\times10^{-16}\) & \(7.62\times10^{-14}\) & \(4.44\times10^{-16}\)
& \(2.81\times10^{-14}\) \\
4 & \(1.42\times10^{-13}\) & \(1.74\times10^{-13}\) &
\(6.66\times10^{-16}\) & \(1.12\times10^{-13}\) & \(4.44\times10^{-16}\)
& \(0\) \\
5 & \(1.82\times10^{-13}\) & \(2.25\times10^{-13}\) &
\(6.66\times10^{-16}\) & \(6.58\times10^{-14}\) & \(4.72\times10^{-16}\)
& \(0\) \\
6 & \(8.91\times10^{-14}\) & \(1.10\times10^{-13}\) &
\(6.66\times10^{-16}\) & \(4.36\times10^{-14}\) & \(5.00\times10^{-16}\)
& \(1.80\times10^{-14}\) \\
7 & \(1.28\times10^{-13}\) & \(1.58\times10^{-13}\) &
\(8.88\times10^{-16}\) & \(4.91\times10^{-14}\) & \(5.00\times10^{-16}\)
& \(2.81\times10^{-14}\) \\
8 & \(1.27\times10^{-13}\) & \(1.55\times10^{-13}\) &
\(1.33\times10^{-15}\) & \(6.42\times10^{-14}\) & \(5.00\times10^{-16}\)
& \(4.60\times10^{-14}\) \\
9 & \(1.70\times10^{-13}\) & \(2.21\times10^{-13}\) &
\(1.67\times10^{-15}\) & \(8.99\times10^{-14}\) & \(4.72\times10^{-16}\)
& \(4.44\times10^{-14}\) \\
\end{longtable}
}

All configurations and trials stayed below the tolerance \(10^{-10}\).
The maximum relabeling-covariance error of the generator was
\(3.33\times10^{-16}\). The null-space component was conserved along the
whole trajectory even for the multi-dimensional kernels (\(N\ge6\)).

\begin{figure}
\centering
\pandocbounded{\includegraphics[keepaspectratio,alt={Conservation and covariance}]{figs/nbody_conservation_errors_v1.png}}
\caption{Conservation and covariance}
\end{figure}

\subsubsection{8.3 Experiment B: generic rank law and degeneracy
search}\label{experiment-b-generic-rank-law-and-degeneracy-search}

For \(N=3\) to \(12\) with 256 uniform random-phase trials each (2560
trials in total),

\[
\operatorname{rank}K=2\min\!\left(N,\left\lfloor\frac{M}{2}\right\rfloor\right)
\]

held in every trial.

{\def\LTcaptype{none} % do not increment counter
\begin{longtable}[]{@{}
  >{\raggedleft\arraybackslash}p{(\linewidth - 12\tabcolsep) * \real{0.1429}}
  >{\raggedleft\arraybackslash}p{(\linewidth - 12\tabcolsep) * \real{0.1429}}
  >{\raggedleft\arraybackslash}p{(\linewidth - 12\tabcolsep) * \real{0.1429}}
  >{\raggedleft\arraybackslash}p{(\linewidth - 12\tabcolsep) * \real{0.1429}}
  >{\raggedleft\arraybackslash}p{(\linewidth - 12\tabcolsep) * \real{0.1429}}
  >{\raggedleft\arraybackslash}p{(\linewidth - 12\tabcolsep) * \real{0.1429}}
  >{\raggedleft\arraybackslash}p{(\linewidth - 12\tabcolsep) * \real{0.1429}}@{}}
\toprule\noalign{}
\begin{minipage}[b]{\linewidth}\raggedleft
\(N\)
\end{minipage} & \begin{minipage}[b]{\linewidth}\raggedleft
\(M\)
\end{minipage} & \begin{minipage}[b]{\linewidth}\raggedleft
Expected rank
\end{minipage} & \begin{minipage}[b]{\linewidth}\raggedleft
Generic observed
\end{minipage} & \begin{minipage}[b]{\linewidth}\raggedleft
Equal phases
\end{minipage} & \begin{minipage}[b]{\linewidth}\raggedleft
\(\{0,\pi\}\) phases
\end{minipage} & \begin{minipage}[b]{\linewidth}\raggedleft
Two-valued \(\{0,0.5\}\) phases
\end{minipage} \\
\midrule\noalign{}
\endhead
\bottomrule\noalign{}
\endlastfoot
3 & 3 & 2 & 2 & 0 & 0 & 2 \\
4 & 6 & 6 & 6 & 0 & 0 & 6 \\
5 & 10 & 10 & 10 & 0 & 0 & \textbf{8} \\
6 & 15 & 12 & 12 & 0 & 0 & 12 \\
7 & 21 & 14 & 14 & 0 & 0 & 14 \\
8 & 28 & 16 & 16 & 0 & 0 & 16 \\
9 & 36 & 18 & 18 & 0 & 0 & 18 \\
10 & 45 & 20 & 20 & 0 & 0 & 20 \\
11 & 55 & 22 & 22 & 0 & 0 & 22 \\
12 & 66 & 24 & 24 & 0 & 0 & 24 \\
\end{longtable}
}

\begin{figure}
\centering
\pandocbounded{\includegraphics[keepaspectratio,alt={Three-layer scaling}]{figs/three_layer_scaling_v1.png}}
\caption{Three-layer scaling}
\end{figure}

\subsubsection{8.4 Exact rational
witnesses}\label{exact-rational-witnesses}

For the witness configuration \(t_e=e\) of Theorem 2, the rank of
\(\widetilde K\) was computed exactly by Gaussian elimination in
fraction arithmetic. For all \(N=3\) to \(12\),
\(\operatorname{rank}\widetilde K=2\min(N,\lfloor M/2\rfloor)\) held
exactly. The computation involves no floating point.

\subsubsection{8.5 Experiment C: demonstration of normal
non-uniqueness}\label{experiment-c-demonstration-of-normal-non-uniqueness}

For display dimensions \(d=3,4,5\) (32 trials each), the
normal-candidate space of two independent directions \(u,v\) was
constructed, and for 64 samples of orthogonal (gauge) transformations
inside the candidate space, the projector drift and the maximum angle
between candidate normals were measured.

{\def\LTcaptype{none} % do not increment counter
\begin{longtable}[]{@{}
  >{\raggedleft\arraybackslash}p{(\linewidth - 8\tabcolsep) * \real{0.2105}}
  >{\raggedleft\arraybackslash}p{(\linewidth - 8\tabcolsep) * \real{0.2105}}
  >{\raggedleft\arraybackslash}p{(\linewidth - 8\tabcolsep) * \real{0.2105}}
  >{\raggedleft\arraybackslash}p{(\linewidth - 8\tabcolsep) * \real{0.2105}}
  >{\raggedright\arraybackslash}p{(\linewidth - 8\tabcolsep) * \real{0.1579}}@{}}
\toprule\noalign{}
\begin{minipage}[b]{\linewidth}\raggedleft
\(d\)
\end{minipage} & \begin{minipage}[b]{\linewidth}\raggedleft
Candidate dim \(d-2\)
\end{minipage} & \begin{minipage}[b]{\linewidth}\raggedleft
Projector gauge drift
\end{minipage} & \begin{minipage}[b]{\linewidth}\raggedleft
Max angle between candidate normals
\end{minipage} & \begin{minipage}[b]{\linewidth}\raggedright
Unique up to sign
\end{minipage} \\
\midrule\noalign{}
\endhead
\bottomrule\noalign{}
\endlastfoot
3 & 1 & \(4.44\times10^{-16}\) & \(0.0000^\circ\) & holds \\
4 & 2 & \(9.99\times10^{-16}\) & \(89.9986^\circ\) & fails \\
5 & 3 & \(1.11\times10^{-15}\) & \(89.9998^\circ\) & fails \\
\end{longtable}
}

For \(d=3\) the gauge transformation does not move the candidate normal;
for \(d\ge4\) the candidate normal moves nearly to orthogonality while
the projector stays invariant. This is the numerical counterpart of
Theorem 3. \(d=4,5\) were not adopted as spaces; they are comparison
targets demonstrating non-uniqueness.

\begin{figure}
\centering
\pandocbounded{\includegraphics[keepaspectratio,alt={Normal gauge freedom}]{figs/normal_gauge_freedom_v1.png}}
\caption{Normal gauge freedom}
\end{figure}

\subsubsection{8.6 Experiment D: ABCD subsystems of N =
5}\label{experiment-d-abcd-subsystems-of-n-5}

For 8 trials at \(N=5\), each of the five ABCD subsystems (40 cases in
total) had its local generator constructed from the six subsystem
relational waves with the same coupling rule. All 40 cases gave local
generator rank 6 with three rotation planes --- the same structure as
the standalone ABCD experiment {[}4{]}.

Moreover, the unnormalized local generator agreed exactly (maximum
difference \(0\)) with the corresponding principal submatrix of the
\(N\)-body master generator. Because endpoint-sharing adjacency closes
within the subsystem, the two agree identically. The two implementations
of subsystem readout --- principal submatrix or local reconstruction ---
are therefore the same object, differing only by the constant of
spectral normalization.

\subsubsection{8.7 Experiment E: residual subspace of multi-dimensional
kernels}\label{experiment-e-residual-subspace-of-multi-dimensional-kernels}

For \(N=6\) to \(9\) (8 trials each), the kernel projector idempotence
error was at most \(1.11\times10^{-15}\) and the annihilation error
\(\lVert KP_0\rVert\) at most \(1.81\times10^{-16}\). Kernel dimensions
equaled \(M-2N\) in every trial.

For two independently chosen orthonormal kernel bases, the projector
difference was at most \(6.66\times10^{-16}\) --- the projector is
unique. The angle between the first basis directions was at least
\(23.5^\circ\) --- the internal basis is not unique. This is a direct
demonstration that the multi-dimensional kernel is a residual subspace
whose projector and squared quantities are readable while its individual
directions are not selectable.

\begin{figure}
\centering
\pandocbounded{\includegraphics[keepaspectratio,alt={Residual subspace growth}]{figs/kernel_dimension_growth_v1.png}}
\caption{Residual subspace growth}
\end{figure}

\begin{center}\rule{0.5\linewidth}{0.5pt}\end{center}

\subsection{9. Scope of Claims}\label{scope-of-claims}

\subsubsection{9.1 Mathematical consequences
proven}\label{mathematical-consequences-proven}

\begin{enumerate}
\def\labelenumi{\arabic{enumi}.}
\tightlist
\item
  Vertex decomposition
  \(\widetilde K=\sum_k(c_ks_k^{\mathsf T}-s_kc_k^{\mathsf T})\)
  (Theorem 1).
\item
  The rank bound \(\operatorname{rank}K\le2\min(N,\lfloor M/2\rfloor)\)
  for all \(N\) and phases (Corollary 1).
\item
  Generic equality for \(3\le N\le12\) (Theorem 2; exact rational
  witnesses + measure-zero argument).
\item
  Sign-free normal uniqueness only for \(d=3\) (Theorem 3).
\end{enumerate}

\subsubsection{9.2 Numerical facts
confirmed}\label{numerical-facts-confirmed}

\begin{enumerate}
\def\labelenumi{\arabic{enumi}.}
\tightlist
\item
  For \(N=3\)--\(9\), 32 trials × 720 steps, closure, absolute-square
  sum, orthogonality, and relabeling covariance were maintained within
  tolerance \(10^{-10}\).
\item
  The generic rank law held in all 2560 trials for \(N=3\)--\(12\).
\item
  Two-valued phases \(\{0,0.5\}\) lowered the rank only at the critical
  configuration \(N=5\) (\(M=2N\)).
\item
  All ABCD subsystems of \(N=5\) carried the three-rotation-plane
  structure, and local generators agreed exactly with principal
  submatrices of the master generator.
\item
  For the multi-dimensional kernels at \(N\ge6\), projector uniqueness
  and internal-basis non-uniqueness held simultaneously.
\end{enumerate}

\subsubsection{9.3 Spatial-direction interpretations
adopted}\label{spatial-direction-interpretations-adopted}

\begin{enumerate}
\def\labelenumi{\arabic{enumi}.}
\tightlist
\item
  Up to three distinguishable rotation planes are read as the XYZ
  directions (inherited from {[}4{]}).
\item
  Rotation planes exceeding the uniqueness condition are internal phase
  modes.
\item
  Multi-dimensional kernels are residual subspaces without a unique
  basis.
\end{enumerate}

\subsubsection{9.4 What this paper does not
fix}\label{what-this-paper-does-not-fix}

\begin{enumerate}
\def\labelenumi{\arabic{enumi}.}
\tightlist
\item
  Generic equality for all \(N\) (stated explicitly as a new
  hypothesis).
\item
  The physical axes corresponding to residual subspaces and internal
  phase modes.
\item
  The mechanism by which three-direction structures of different
  subsystems become phase-locked to the same XYZ under dynamics.
\item
  Whether the three-direction ceiling persists when higher-order
  relations beyond pairwise are added.
\item
  Derivation of the phase-difference sine generator itself from the
  axioms.
\item
  Properties of the sequentially reconstructed generator (dynamics
  rebuilding the generator from current phases at each step). This paper
  concerns the fixed generator.
\end{enumerate}

\begin{center}\rule{0.5\linewidth}{0.5pt}\end{center}

\subsection{10. Reproducibility}\label{reproducibility}

The numerical experiments of this paper are based on the following
programs and outputs.

\subsubsection{Programs}\label{programs}

\begin{itemize}
\tightlist
\item
  \texttt{dimension\_generation\_structure/contrast\_experiments\_v1/run\_nbody\_rank\_saturation\_preliminary\_v1.py}
  (Experiments A--E)
\item
  \texttt{dimension\_generation\_structure/contrast\_experiments\_v1/run\_nbody\_exact\_rational\_witness\_v1.py}
  (exact rational witnesses)
\item
  \texttt{dimension\_generation\_structure/contrast\_experiments\_v1/make\_nbody\_rank\_saturation\_figures\_v1.py}
  (figures)
\end{itemize}

\subsubsection{Main outputs}\label{main-outputs}

\begin{itemize}
\tightlist
\item
  \texttt{nbody\_rank\_saturation\_preliminary\_result\_v1.json}
\item
  \texttt{experiment\_a\_trials\_v1.csv} /
  \texttt{experiment\_a\_summary\_v1.csv}
\item
  \texttt{experiment\_b\_rank\_law\_v1.csv}
\item
  \texttt{experiment\_c\_normal\_uniqueness\_v1.csv}
\item
  \texttt{experiment\_d\_subsystems\_v1.csv}
\item
  \texttt{experiment\_e\_kernel\_v1.csv}
\item
  \texttt{nbody\_exact\_rational\_witness\_v1.json}
\item
  \texttt{nbody\_rank\_saturation\_preliminary\_report\_v1.md}
\end{itemize}

These are stored in

\texttt{dimension\_generation\_structure/figs/}.

\begin{center}\rule{0.5\linewidth}{0.5pt}\end{center}

\subsection{11. Final Conclusions}\label{final-conclusions}

The \(N\)-body complete pairwise relational-wave system was constructed
without placing background spatial axes first, on the basis of the
complete pairwise relational waves and squared closure of Axiom 16, with
the phase-difference sine generator as a conservative working hypothesis
{[}1{]}.

The number of internal relational waves grows quadratically as
\(M_N=\binom{N}{2}\). The generator, however, decomposes exactly into a
per-vertex sum of rank-2 antisymmetric matrices built from the incidence
rows, giving

\[
\operatorname{rank}K\le2\min\!\left(N,\left\lfloor\frac{M_N}{2}\right\rfloor\right)
\]

for all \(N\). For \(3\le N\le12\), exact rational witness
configurations combined with the measure-zero property of zero sets of
real-analytic functions prove that the bound is an equality in general
position. Rotation modes grow at most linearly.

In the direction readout that reads a normal from a two-relation plane,
the normal-candidate space has dimension \(d-2\), and sign-free
uniqueness holds only for \(d=3\). Even when distinguishable rotation
planes increase beyond three, the excess planes cannot be adopted as
spatial directions and remain internal phase modes. The
multi-dimensional kernel reappearing at \(N\ge6\) is a residual subspace
with a unique projector but no selectable internal basis.

The conclusions of this paper are therefore:

\[
\boxed{
\begin{aligned}
&\text{internal relational waves grow as }O(N^2);\\
&\text{internal rotation modes do not exceed }O(N);\\
&\text{under the present two-body phase-relation readout rule, uniquely readable spatial directions saturate at three.}
\end{aligned}
}
\]

Which physical axes the residual subspaces and internal phase modes
correspond to is outside the scope of this paper.

\begin{center}\rule{0.5\linewidth}{0.5pt}\end{center}

\subsection{References}\label{references}

\subsubsection{Self-citations}\label{self-citations}

\begin{enumerate}
\def\labelenumi{\arabic{enumi}.}
\tightlist
\item
  Noriaki Kihara, ``Anonymous Equal-Amplitude Composite-Wave Model:
  Basic Axiom System v5,'' Zenodo, 2026. Version DOI:
  \href{https://doi.org/10.5281/zenodo.21465429}{10.5281/zenodo.21465429},
  Concept DOI:
  \href{https://doi.org/10.5281/zenodo.21315735}{10.5281/zenodo.21315735}.
\item
  Noriaki Kihara, ``Harmonic Readout and c=1 Area-Sweep Preliminary
  Experiments in the Closed AB Two-Body Phase System, v4,'' Zenodo,
  2026. Version DOI:
  \href{https://doi.org/10.5281/zenodo.21374317}{10.5281/zenodo.21374317},
  Concept DOI:
  \href{https://doi.org/10.5281/zenodo.21318696}{10.5281/zenodo.21318696}.
\item
  Noriaki Kihara, ``Inverse-Square Law from the Future-Phase-Position
  Acceleration Map and Harmonic Closure in the Closed AB Two-Body Phase
  System,'' Zenodo, 2026. Version DOI:
  \href{https://doi.org/10.5281/zenodo.21441082}{10.5281/zenodo.21441082},
  Concept DOI:
  \href{https://doi.org/10.5281/zenodo.21441081}{10.5281/zenodo.21441081}.
\item
  Noriaki Kihara, ``Three XYZ Directions Read from Complete Pairwise
  Relational Waves: Growth of Internal Relational Directions and
  Saturation of Spatial-Direction Readout in Closed AB, ABC, and ABCD
  Systems,'' Zenodo, 2026. Version DOI:
  \href{https://doi.org/10.5281/zenodo.21454790}{10.5281/zenodo.21454790},
  Concept DOI:
  \href{https://doi.org/10.5281/zenodo.21454789}{10.5281/zenodo.21454789}.
\end{enumerate}

\subsubsection{External references}\label{external-references}

\begin{enumerate}
\def\labelenumi{\arabic{enumi}.}
\setcounter{enumi}{4}
\tightlist
\item
  Hassler Whitney, ``Congruent Graphs and the Connectivity of Graphs,''
  \emph{American Journal of Mathematics}, 54(1), 150--168, 1932. DOI:
  \href{https://doi.org/10.2307/2371086}{10.2307/2371086}.
\item
  Milan Vujivčić, Fedor Herbut, and Gradimir Vujivčić, ``Canonical Form
  for Matrices Under Unitary Congruence Transformations. I:
  Conjugate-Normal Matrices,'' \emph{SIAM Journal on Applied
  Mathematics}, 23(2), 225--238, 1972. DOI:
  \href{https://doi.org/10.1137/0123025}{10.1137/0123025}.
\item
  Fasma Diele, Luciano Lopez, and R. Peluso, ``The Cayley Transform in
  the Numerical Solution of Unitary Differential Systems,''
  \emph{Advances in Computational Mathematics}, 8(4), 317--334, 1998.
  DOI:
  \href{https://doi.org/10.1023/A:1018908700358}{10.1023/A:1018908700358}.
\item
  Cyriel van Nuffelen, ``On the incidence matrix of a graph,''
  \emph{IEEE Transactions on Circuits and Systems}, 23(9), 572, 1976.
  DOI:
  \href{https://doi.org/10.1109/TCS.1976.1084251}{10.1109/TCS.1976.1084251}.
\item
  Boris S. Mityagin, ``The Zero Set of a Real Analytic Function,''
  \emph{Mathematical Notes}, 107(3), 529--530, 2020. DOI:
  \href{https://doi.org/10.1134/S0001434620030189}{10.1134/S0001434620030189}.
\end{enumerate}

\end{document}
